\chapter{自由场与多粒子态}\label{chap:qft-foundations}

\cref{chap:classical-gauge-response}已经确定经典场的耦合、守恒电流与因果响应；场本身仍是给定的动力学变量。量子理论还必须容纳相对论相互作用中的粒子产生与湮灭。若预先只选一个一电子 Hilbert 空间，那么一旦过程允许出现电子--正电子对或真实光子，这个空间就不再对时间演化封闭。量子场论首先要解决的正是这个状态空间问题。

经典场可以分解成许多正常模，每个模在自由极限下都像一个谐振子；量子化这些模以后，产生、湮灭算符改变模的占据数，把零粒子、一粒子和多粒子状态统一放进 Fock 空间。连续谱归一化必须与场的对易关系和能量同时相容，局域场还必须满足类空分离处的微观因果性。自由状态空间建立以后，\cref{chap:interacting-field-propagation}才能在其中定义传播与相互作用。

实标量场在这里承担最小模型的角色：没有旋量指标，也没有规范冗余，因此模式、占据数和局域对易关系可以先单独显现出来。它描述自旋 $0$ 场；自旋 $1/2$ 电子仍由 Dirac 算符描述，光子仍由具有规范冗余的 Maxwell 场描述。相应场的量子化会在这些一般对象建立以后分别展开。

\section{自由场量子化与多粒子状态空间}

经典场的正常模已经像一组彼此独立的谐振子。把这些模式量子化后，每个模的占据数成为离散量子数；对带电场，正、负频分支进一步对应粒子和反粒子。这样得到的 Fock 空间同时容纳不同粒子数，传播子和散射外态才可以在同一 Hilbert 空间中定义。

连续动量极限要求场算符、单粒子态与 Lorentz 不变相空间采用彼此兼容的归一化；否则二点传播核、从相关函数提取的单粒子外态以及散射相空间中的能量因子将失去共同来源。实标量场给出这套归一化的最小模型，复标量场则在相同归一化下进一步区分粒子与反粒子的两套独立产生、湮灭算符。

\subsection{量子场的必要性}

非相对论单粒子量子力学预先选定一个一粒子 Hilbert 空间，随后研究这个空间内的波函数演化。相对论把能量与质量联系起来，因此足够集中的能量可以转化为新的粒子自由度。设质量为 $m$ 的波包被局域在尺度 $L$ 内，Heisenberg 关系给出典型动量不确定度
\begin{equation*}
 \Delta p\sim \frac{\hbar}{L}.
\end{equation*}
当 $L$ 缩短到使相对论动能尺度 $c\Delta p$ 与静能 $mc^2$ 同阶时，固定粒子数描述失去封闭性：
\begin{equation}
 L\sim \lambda_C,
 \qquad
 \lambda_C\equiv\frac{\hbar}{mc}.
 \label{eq:qft-compton-scale}
\end{equation}
若过程允许提供约 $2mc^2$ 的能量，则粒子--反粒子对可以进入动力学。阈值的精确数值取决于具体过程，而 $\lambda_C$ 给出单粒子局域化概念开始失效的基本长度尺度。

电子的相对论动力学应直接从 Dirac 方程理解。它的一阶时间演化具有正定密度 $\psi^\dagger\psi$，但自由解仍同时包含正、负频率分支；量子化以后这两支分别组织成电子湮灭算符和正电子产生算符。更重要的是，局域相互作用能够连接不同粒子数扇区，所以固定粒子数 Hilbert 空间仍不是相对论相互作用理论的封闭状态空间。Fock 空间因此来自粒子产生与湮灭的动力学需要，而不是由某一种二阶波动方程的缺陷“补救”出来。

场论同时实现了另一个要求：局域性。粒子间的“远距势能”可改写成局域场与源的耦合；扰动先改变附近的场，再按场方程传播。经典场论已经用局域场把相互作用的传播写成微分方程。量子场论保留这一局域结构，但把场本身提升为算符；不同场在既定时空对称性下属于不同表示，其稳定单量子谱由此表现为不同粒子种类。

经典场 $\varphi(\vb{x},t)$ 中的 $\vb{x}$ 是自由度的标签，而不是某个粒子的坐标。量子化后，$\hat\varphi(\vb{x},t)$ 对每个空间点都给出一个算符，但这种“逐点算符”在连续极限中通常含有 $\delta$ 函数型奇性；严格可作用在态上的对象是把它与光滑、有限支撑的测试函数积分后的涂抹场。因此量子场严格说是\emph{算符值分布}：它像数学物理方法中的 $\delta$ 分布一样，通过积分作用定义，而不是要求每个点上的算符都单独有界。对场作正常模展开后，相应的产生算符会构造具有确定能量和动量的单量子态。于是“场”是动力学变量，“粒子”是场的谱结构。因此场的运动方程与对称性先确定允许的谱结构，粒子种类再由该谱中的稳定或共振激发识别。

\subsection{标量场作为最小高斯模型}

电子最终要量子化的是 Dirac 场，但正则量子化、Fock 空间、二点函数、Dyson 展开和 Wick 配对的共同数学结构与自旋无关。为了先把这些结构算清楚，这里暂时取一个自旋 $0$ 的实标量场；它只承担最小高斯模型的角色，后续电子结果由 Dirac 场重新建立。其自由拉氏量取
\begin{equation}
 \Lag_0
 =\frac12\left[
 \frac{1}{c^2}(\partial_t\varphi)^2
 -(\vb{\nabla}\varphi)^2
 -\frac{m^2c^2}{\hbar^2}\varphi^2
 \right],
 \label{eq:kg-lagrangian-r2}
\end{equation}
则 Euler--Lagrange 方程为
\begin{equation}
 \left(\Box+\frac{m^2c^2}{\hbar^2}\right)\varphi=0,
 \qquad
 \Box\equiv\partial_\mu\partial^\mu
 =\frac1{c^2}\partial_t^2-\nabla^2.
 \label{eq:kg-eom-r2}
\end{equation}
共轭动量与哈密顿量密度分别为
\begin{align}
 \pi(\vb{x},t)
 &=\frac{\partial\Lag_0}{\partial(\partial_t\varphi)}
 =\frac1{c^2}\partial_t\varphi,
 \label{eq:kg-pi-r2}\\
 \Ham_0
 &=\pi\partial_t\varphi-\Lag_0
 =\frac{c^2}{2}\pi^2
 +\frac12(\nabla\varphi)^2
 +\frac{m^2c^2}{2\hbar^2}\varphi^2.
 \label{eq:kg-H-density-r2}
\end{align}
这已经是一个具有连续无穷多个自由度的哈密顿系统。

\subsection{正则量子化}

经典自由场已经被分解为彼此独立的正常模；量子化所要做的，是把每个正常模的经典正则变量提升为算符，并要求它们满足与经典 Poisson 括号对应的正则对易关系。先把系统放在体积 $V=L^3$ 的周期盒中，允许动量离散取值
\begin{equation*}
 \vb{p}=\frac{2\pi\hbar}{L}\vb{n},
 \qquad \vb{n}\in\mathbb Z^3.
\end{equation*}
实标量场可写成
\begin{equation*}
 \varphi(\vb{x},t)
 =\frac1{\sqrt V}\sum_{\vb{p}}q_{\vb{p}}(t)
 \ee^{\ii\vb{p}\cdot\vb{x}/\hbar},
 \qquad q_{-\vb{p}}=q_{\vb{p}}^*.
\end{equation*}
代入 Klein--Gordon 方程 \eqnrefs{eq:kg-eom-r2} 后，各动量模彼此解耦，并分别满足
\begin{equation}
 \ddot q_{\vb{p}}+\omega_{\vb{p}}^2q_{\vb{p}}=0,
 \qquad
 \omega_{\vb{p}}\equiv\frac{E_{\vb{p}}}{\hbar},
 \qquad
 E_{\vb{p}}=\sqrt{m^2c^4+c^2\vb{p}^2}.
 \label{eq:kg-mode-frequency-r2}
\end{equation}
因此，自由场不是一种全新的力学系统，而是无限多个频率不同的谐振子在连续极限下的集合；量子化以后，每个正常模自然对应一组产生与湮灭算符。

对经过正则缩放的单个谐振子，取
\begin{equation*}
 H_{\rm osc}=\frac{P^2}{2}+\frac{\omega^2Q^2}{2},
 \qquad [Q,P]=\ii\hbar,
\end{equation*}
并定义
\begin{equation}
 a=\sqrt{\frac{\omega}{2\hbar}}Q
 +\frac{\ii}{\sqrt{2\hbar\omega}}P,
 \qquad
 a^\dagger=\sqrt{\frac{\omega}{2\hbar}}Q
 -\frac{\ii}{\sqrt{2\hbar\omega}}P.
 \label{eq:sho-ladder-r3}
\end{equation}
这一定义把正则对易关系等价地写成
\begin{equation}
 [a,a^\dagger]=1,
 \qquad
 H_{\rm osc}=\hbar\omega\left(a^\dagger a+\frac12\right).
 \label{eq:sho-H-r3}
\end{equation}
于是 $a^\dagger$ 与 $a$ 分别使能量增加或减少一个 $\hbar\omega$；若能谱有下界，真空满足 $a|0\rangle=0$，全部激发可由 $(a^\dagger)^n|0\rangle$ 构造。场论只是在每个连续动量标签 $\vb{p}$ 上复制这套结构。


\begin{exbox}[把一个场模真的当作谐振子来算]
为使“一个场模就是一个量子谐振子”成为可数的离散问题，把质量 $m>0$ 的实标量场先放在边长为 $L$ 的周期盒中。最低的零动量模满足 $\vb{p}=\mathbf0$，因此
\begin{equation*}
 E_{\mathbf0}=mc^2,
 \qquad
 \omega_{\mathbf0}=\frac{mc^2}{\hbar}.
\end{equation*}
计算零动量模的零点能和前两个激发能，并与第一个非零动量模比较。
\textbf{解。} 量子化以后，这个模的哈密顿量就是
\begin{equation*}
 H_{\mathbf0}=mc^2\left(a_{\mathbf0}^\dagger a_{\mathbf0}+\frac12\right).
\end{equation*}
这一模的 $mc^2/2$ 是所有占据数态共有的零点常数。为了只比较相对于真空的激发能，先在这个例题中取
\begin{equation*}
 H_{\mathbf0}^{\rm exc}
 \equiv H_{\mathbf0}-\langle0|H_{\mathbf0}|0\rangle
 =mc^2a_{\mathbf0}^\dagger a_{\mathbf0}.
\end{equation*}
于是，真空、一个零动量量子和两个零动量量子依次满足
\begin{align*}
 H_{\mathbf0}^{\rm exc}|0\rangle&=0,\\
 H_{\mathbf0}^{\rm exc}a_{\mathbf0}^\dagger|0\rangle
 &=mc^2a_{\mathbf0}^\dagger|0\rangle,\\
 H_{\mathbf0}^{\rm exc}\frac{(a_{\mathbf0}^\dagger)^2}{\sqrt{2!}}|0\rangle
 &=2mc^2\frac{(a_{\mathbf0}^\dagger)^2}{\sqrt{2!}}|0\rangle.
\end{align*}
这就是“静止质量”在场量子化中的最直接出现方式：$mc^2$ 不是另外塞进粒子模型的能量，而是零动量正常模的单量子激发能。

再看最小的非零动量，例如 $\vb{p}=(2\pi\hbar/L)\vb{e}_x$。它的一次激发能变为
\begin{equation*}
 E_{1x}=\sqrt{m^2c^4+c^2\left(\frac{2\pi\hbar}{L}\right)^2}.
\end{equation*}
于是同一套升降算符代数自动同时生成不同动量、不同粒子数的态。取 $L\to\infty$ 后，离散标签 $\vb{p}$ 变成连续变量，这给出连续 Fock 空间的极限来源，而不是一次新的物理假设。
\end{exbox}

\begin{exbox}[两个动量模的 Fock 态]
“一个量子”与“两个量子”在场论里不是不同种类的波函数，而是同一组产生算符在 Fock 空间中作用不同次数。为把这句话算成一个可以直接检查的结果，仍在有限周期盒中取两个不同动量模 $\vb{p}\neq\vb{q}$，并定义归一化占据数态
\begin{equation*}
 |n_{\vb{p}},n_{\vb{q}}\rangle
 =\frac{(a_{\vb{p}}^\dagger)^{n_{\vb{p}}}}{\sqrt{n_{\vb{p}}!}}
  \frac{(a_{\vb{q}}^\dagger)^{n_{\vb{q}}}}{\sqrt{n_{\vb{q}}!}}|0\rangle .
\end{equation*}
求这些态的总能量、总动量，以及吸收和发射一个量子的阶梯系数。
\textbf{解。} 两个不同模彼此对易。把自由哈密顿量的共同真空常数减去，并记所得激发哈密顿量为 $H_0^{\rm exc}$，则
\begin{align*}
 H_0^{\rm exc}|n_{\vb{p}},n_{\vb{q}}\rangle
 &=\bigl(n_{\vb{p}}E_{\vb{p}}+n_{\vb{q}}E_{\vb{q}}\bigr)
 |n_{\vb{p}},n_{\vb{q}}\rangle,\\
 \vb{P}\,|n_{\vb{p}},n_{\vb{q}}\rangle
 &=\bigl(n_{\vb{p}}\vb{p}+n_{\vb{q}}\vb{q}\bigr)
 |n_{\vb{p}},n_{\vb{q}}\rangle.
\end{align*}
例如 $|2_{\vb{p}},1_{\vb{q}}\rangle$ 的能量和总动量分别是
\begin{equation*}
 2E_{\vb{p}}+E_{\vb{q}},
 \qquad
 2\vb{p}+\vb{q}.
\end{equation*}
这说明 Fock 态中的“粒子数”就是各正常模的占据数，而能量和动量是各量子的逐个相加。

再看同一个模上的产生、湮灭算符。由 $[a_{\vb{p}},a_{\vb{p}}^\dagger]=1$，逐次把 $a_{\vb{p}}$ 移过 $(a_{\vb{p}}^\dagger)^n$，得到
\begin{align*}
 a_{\vb{p}}(a_{\vb{p}}^\dagger)^n|0\rangle
 &=n(a_{\vb{p}}^\dagger)^{n-1}|0\rangle,\\
 a_{\vb{p}}|n_{\vb{p}}\rangle
 &=\sqrt{n_{\vb{p}}}\,|n_{\vb{p}}-1\rangle,\\
 a_{\vb{p}}^\dagger|n_{\vb{p}}\rangle
 &=\sqrt{n_{\vb{p}}+1}\,|n_{\vb{p}}+1\rangle.
\end{align*}
所以相互作用哈密顿量中若出现一个 $a_{\vb{p}}$ 或 $a_{\vb{p}}^\dagger$，吸收与发射振幅分别带 $\sqrt n$ 与 $\sqrt{n+1}$。自由电子与量子光场相互作用中的玻色增强因子也来自同一个 Fock 阶梯代数。
\end{exbox}

回到场变量，正则量子化要求等时算符满足
\begin{align}
 [\hat\varphi(\vb{x},t),\hat\varphi(\vb{y},t)]&=0,\\
 [\hat\pi(\vb{x},t),\hat\pi(\vb{y},t)]&=0,\\
 [\hat\varphi(\vb{x},t),\hat\pi(\vb{y},t)]
 &=\ii\hbar\delta^{(3)}(\vb{x}-\vb{y}).
 \label{eq:kg-canonical-r2}
\end{align}
连续动量表象中取产生、湮灭算符的归一化
\begin{equation}
 [\hat a(\vb{p}),\hat a^\dagger(\vb{q})]
 =(2\pi\hbar)^3\delta^{(3)}(\vb{p}-\vb{q}).
 \label{eq:scalar-aa-r2}
\end{equation}

\begin{proposition}[正则场的模式 CCR]
设自由实标量场采用本节的 Fourier 测度及实正模式系数 $A_{\vb{p}}=\hbar c/\sqrt{2E_{\vb{p}}}$，满足等时正则对易关系。则在算符值分布意义下
\[
[a(\vb{p}),a^\dagger(\vb{q})]=(2\pi\hbar)^3\delta^{(3)}(\vb{p}-\vb{q}),
\qquad [a(\vb{p}),a(\vb{q})]=0.
\]
\end{proposition}
\begin{proof}
在 $t=0$ 的等时面，令 $B_{\vb{p}}=\sqrt{E_{\vb{p}}/(2c^2)}$。模式展开的 Fourier 反演给出
\[
a(\vb{p})=\frac12\int\dd^3x\,\ee^{-\ii\vb{p}\cdot\vb{x}/\hbar}
\left(\frac{\varphi(\vb{x})}{A_{\vb{p}}}+\frac{\ii\pi(\vb{x})}{B_{\vb{p}}}\right).
\]
将此式及其伴随代入对易子，场--场与动量--动量项为零，而 $[\varphi(\vb{x}),\pi(\vb{y})]=\ii\hbar\delta^{(3)}(\vb{x}-\vb{y})$ 给
\begin{align*}
[a(\vb{p}),a^\dagger(\vb{q})]
&=\frac\hbar4\left(\frac1{A_{\vb{p}}B_{\vb{q}}}+
\frac1{B_{\vb{p}}A_{\vb{q}}}\right)
\int\dd^3x\,\ee^{\ii(\vb{q}-\vb{p})\cdot\vb{x}/\hbar}\\
&=(2\pi\hbar)^3\delta^{(3)}(\vb{p}-\vb{q}),
\end{align*}
其中在 $\vb{p}=\vb{q}$ 上用了 $A_{\vb{p}}B_{\vb{p}}=\hbar/2$。
两个湮灭算符的对易子则含 $\delta^{(3)}(\vb{p}+\vb{q})$，其系数为
$\hbar(-1/(A_{\vb{p}}B_{\vb{q}})+1/(B_{\vb{p}}A_{\vb{q}}))/4$，由 $E_{-\vb{p}}=E_{\vb{p}}$ 为零；产生--产生关系由伴随得到。
\end{proof}

此时平移不变性、实场的厄米条件以及正频率质量壳已经把场展开的形式限制为一对正、负频率平面波；唯一尚未确定的是每个动量模前的归一化系数。这个系数不能任意选择，因为它必须使场的等时对易关系恰好恢复 \eqnrefs{eq:kg-canonical-r2}。比较 Fourier 核后得到


\begin{proposition}[模式系数的正则归一化]
设自由标量场的模式完备，采用正文的 Fourier 测度和等时 CCR，系数取所声明的实正相位约定。则有
\begin{equation}
 C_{\vb{p}}=\frac{\hbar c}{\sqrt{2E_{\vb{p}}}}.
\label{eq:scalar-C-normalization-r16}
\end{equation}
\end{proposition}
\begin{proof}
满足平移不变性和实场厄米条件的一般质量壳展开可写成

\begin{equation*}
 \hat\varphi(x)=
 \int\frac{\dd^3 p}{(2\pi\hbar)^3}C_{\vb{p}}
 \left[
 \hat a(\vb{p})\ee^{-\ii p\cdot x/\hbar}
 +\hat a^\dagger(\vb{p})\ee^{+\ii p\cdot x/\hbar}
 \right],
 \qquad C_{\vb{p}}=C_{-\vb{p}}>0.
\end{equation*}
由 $\hat\pi=c^{-2}\partial_t\hat\varphi$，
\begin{equation*}
 \hat\pi(x)=
 -\ii\int\frac{\dd^3 p}{(2\pi\hbar)^3}
 \frac{E_{\vb{p}}C_{\vb{p}}}{\hbar c^2}
 \left[
 \hat a(\vb{p})\ee^{-\ii p\cdot x/\hbar}
 -\hat a^\dagger(\vb{p})\ee^{+\ii p\cdot x/\hbar}
 \right].
\end{equation*}
在等时面 $x^0=y^0$ 上，把这两式与 \eqnrefs{eq:scalar-aa-r2} 代入对易子。$aa$ 与 $a^\dagger a^\dagger$ 项为零，只剩
\begin{align*}
 [\hat\varphi(\vb{x},t),\hat\pi(\vb{y},t)]
 &=\ii\int\frac{\dd^3 p}{(2\pi\hbar)^3}
 \frac{E_{\vb{p}}C_{\vb{p}}^{2}}{\hbar c^2}
 \left[
 \ee^{+\ii\vb{p}\cdot(\vb{x}-\vb{y})/\hbar}
 +\ee^{-\ii\vb{p}\cdot(\vb{x}-\vb{y})/\hbar}
 \right]\\
 &=2\ii\int\frac{\dd^3 p}{(2\pi\hbar)^3}
 \frac{E_{\vb{p}}C_{\vb{p}}^{2}}{\hbar c^2}
 \ee^{+\ii\vb{p}\cdot(\vb{x}-\vb{y})/\hbar}.
\end{align*}
第二行在第二个积分中作 $\vb{p}\mapsto-\vb{p}$，并使用 $E_{-\vb{p}}=E_{\vb{p}}$ 与 $C_{-\vb{p}}=C_{\vb{p}}$。另一方面，
\begin{equation*}
 \delta^{(3)}(\vb{x}-\vb{y})
 =\int\frac{\dd^3 p}{(2\pi\hbar)^3}
 \ee^{+\ii\vb{p}\cdot(\vb{x}-\vb{y})/\hbar}.
\end{equation*}
与 \eqnrefs{eq:kg-canonical-r2} 比较同一 Fourier 核的系数，得到
\begin{equation*}
 \frac{2E_{\vb{p}}C_{\vb{p}}^{2}}{\hbar c^2}=\hbar
 \quad\Longrightarrow\quad
 C_{\vb{p}}=\frac{\hbar c}{\sqrt{2E_{\vb{p}}}},
\end{equation*}
即\eqnrefs{eq:scalar-C-normalization-r16}。
\end{proof}

因此自由实标量场的连续模展开为
\begin{equation}
 {
 \hat\varphi(x)=
 \int\frac{\dd^3 p}{(2\pi\hbar)^3}
 \frac{\hbar c}{\sqrt{2E_{\vb{p}}}}
 \left[
 \hat a(\vb{p})\ee^{-\ii p\cdot x/\hbar}
 +\hat a^\dagger(\vb{p})\ee^{+\ii p\cdot x/\hbar}
 \right].}
 \label{eq:scalar-mode-r2}
\end{equation}
其中 $p^\mu=(E_{\vb{p}}/c,\vb{p})$ 位于质量壳 $p^2=m^2c^2$。对\eqnrefs{eq:scalar-mode-r2} 的两个频率分支分别作时间导数，并使用共轭动量定义~\eqnrefs{eq:kg-pi-r2}，得到
\begin{equation*}
 \hat\pi(x)=
 -\ii\int\frac{\dd^3 p}{(2\pi\hbar)^3}
 \sqrt{\frac{E_{\vb{p}}}{2c^2}}
 \left[
 \hat a(\vb{p})\ee^{-\ii p\cdot x/\hbar}
 -\hat a^\dagger(\vb{p})\ee^{+\ii p\cdot x/\hbar}
 \right].
\end{equation*}

产生、湮灭算符的物理意义由时空平移生成元确定，而不只是来自谐振子的类比。把上述模展开代入自由场哈密顿量，质量壳关系使 $aa$ 与 $a^\dagger a^\dagger$ 项相消，留下
\begin{align}
 \hat H_0
 &=\frac12\int_{\vb{p}}E_{\vb{p}}
 \left(
 \hat a_{\vb{p}}\hat a_{\vb{p}}^\dagger
 +\hat a_{\vb{p}}^\dagger\hat a_{\vb{p}}
 \right),\\
 &=\int_{\vb{p}}E_{\vb{p}}
 \left[
 \hat a_{\vb{p}}^\dagger\hat a_{\vb{p}}
 +\frac12(2\pi\hbar)^3\delta^{(3)}(0)
 \right],
 \label{eq:scalar-H-vac-r2}
\end{align}

\begin{proposition}[自由实标量哈密顿量的模式对角化]
采用上述模式展开与 CCR，自由场的二次哈密顿量在有限盒及紫外截断下满足
\[
H_0=\frac12\int_{\vb{p}}E_{\vb{p}}(a_{\vb{p}}a_{\vb{p}}^\dagger+a_{\vb{p}}^\dagger a_{\vb{p}}).
\]
质量壳使 $aa$ 与 $a^\dagger a^\dagger$ 项精确抵消。减去真空常数后，$H_0^{\rm exc}=\int_{\vb{p}}E_{\vb{p}}a_{\vb{p}}^\dagger a_{\vb{p}}$；连续极限按算符值分布与正规序的既定约定理解。
\end{proposition}
\begin{proof}
把

\begin{equation*}
 \hat H_0=\frac12\int\dd^3 x\left[
 c^2\hat\pi^2+(\nabla\hat\varphi)^2
 +\frac{m^2c^2}{\hbar^2}\hat\varphi^2
 \right]
 \end{equation*}
分成时间导数项、空间梯度项与质量项三部分。为缩短中间式，记
\begin{equation*}
 A_{\vb{p}}\equiv\frac{\hbar c}{\sqrt{2E_{\vb{p}}}},
 \qquad \int_{\vb{p}}\equiv\int\frac{\dd^3 p}{(2\pi\hbar)^3}.
\end{equation*}
在固定时间上，\eqnrefs{eq:scalar-mode-r2} 可写成
\begin{align*}
 \hat\varphi(\vb{x})
 &=\int_{\vb{p}}A_{\vb{p}}
 \left(\hat a_{\vb{p}}\ee^{+\ii\vb{p}\cdot\vb{x}/\hbar-\ii E_{\vb{p}}t/\hbar}
 +\hat a_{\vb{p}}^\dagger \ee^{-\ii\vb{p}\cdot\vb{x}/\hbar+\ii E_{\vb{p}}t/\hbar}\right),\\
 \hat\pi(\vb{x})
 &=-\ii\int_{\vb{p}}\sqrt{\frac{E_{\vb{p}}}{2c^2}}
 \left(\hat a_{\vb{p}}\ee^{+\ii\vb{p}\cdot\vb{x}/\hbar-\ii E_{\vb{p}}t/\hbar}
 -\hat a_{\vb{p}}^\dagger \ee^{-\ii\vb{p}\cdot\vb{x}/\hbar+\ii E_{\vb{p}}t/\hbar}\right).
 \end{align*}
先看两个湮灭算符相乘的项。空间积分
\begin{equation*}
 \int\dd^3 x\,
 \ee^{\ii(\vb{p}+\vb{q})\cdot\vb{x}/\hbar}
 =(2\pi\hbar)^3\delta^{(3)}(\vb{p}+\vb{q})
 \end{equation*}
强制 $\vb{q}=-\vb{p}$，从而 $E_{\vb{q}}=E_{\vb{p}}$。三部分分别给出
\begin{align*}
 H_{\pi}^{(aa)}
 &=-\int_{\vb{p}}\frac{E_{\vb{p}}}{4}\,
 \hat a_{\vb{p}}\hat a_{-\vb{p}}
 \ee^{-2\ii E_{\vb{p}}t/\hbar},\\
 H_{\nabla}^{(aa)}
 &=+\int_{\vb{p}}\frac{c^2\vb{p}^2}{4E_{\vb{p}}}\,
 \hat a_{\vb{p}}\hat a_{-\vb{p}}
 \ee^{-2\ii E_{\vb{p}}t/\hbar},\\
 H_{m}^{(aa)}
 &=+\int_{\vb{p}}\frac{m^2c^4}{4E_{\vb{p}}}\,
 \hat a_{\vb{p}}\hat a_{-\vb{p}}
 \ee^{-2\ii E_{\vb{p}}t/\hbar}.
 \end{align*}
所以总系数只需一行：
\begin{equation*}
 -\frac{E_{\vb{p}}}4+\frac{c^2\vb{p}^2+m^2c^4}{4E_{\vb{p}}}
 =\frac{-E_{\vb{p}}^2+c^2\vb{p}^2+m^2c^4}{4E_{\vb{p}}}=0.
\end{equation*}
这里仅使用自由质量壳 $E_{\vb{p}}^2=c^2\vb{p}^2+m^2c^4$。$a^\dagger a^\dagger$ 项是其厄米共轭，因此同样消失。这些项的消失源于自由标量二次作用量与其质量壳频率的精确匹配；不存在额外的粒子数守恒删项。

再看一个产生算符与一个湮灭算符相乘的交叉项。空间积分给 $\delta^{(3)}(\vb{p}-\vb{q})$，三部分系数直接合并为
\begin{equation*}
 \frac{E_{\vb{p}}}4+\frac{c^2\vb{p}^2+m^2c^4}{4E_{\vb{p}}}
 =\frac{E_{\vb{p}}^2+c^2\vb{p}^2+m^2c^4}{4E_{\vb{p}}}
 =\frac{E_{\vb{p}}}2.
\end{equation*}
因此
\begin{align*}
 \hat H_0
 &=\frac12\int_{\vb{p}}E_{\vb{p}}
 \left(
 \hat a_{\vb{p}}\hat a_{\vb{p}}^\dagger
 +\hat a_{\vb{p}}^\dagger\hat a_{\vb{p}}
 \right),\\
 &=\int_{\vb{p}}E_{\vb{p}}
 \left[
 \hat a_{\vb{p}}^\dagger\hat a_{\vb{p}}
 +\frac12(2\pi\hbar)^3\delta^{(3)}(0)
 \right],
 \end{align*}
其中第二行用了 \eqnrefs{eq:scalar-aa-r2}。这样单个谐振子的 $\hbar\omega(a^\dagger a+1/2)$ 与连续场的结果逐项对应，因为 $\hbar\omega_{\vb{p}}=E_{\vb{p}}$。
\end{proof}
其中
\begin{equation*}
 \int_{\vb{p}}\equiv\int\frac{\dd^3 p}{(2\pi\hbar)^3}.
\end{equation*}
第二项是所有正常模零点能的连续和。若当前只研究相对于这个真空的粒子激发，可以先把共同的真空常数从能量零点中分离出来，并在普通正文中定义激发哈密顿量
\begin{equation}
 {
 \hat H_0^{\rm exc}
 \equiv \hat H_0-\langle0|\hat H_0|0\rangle
 =\int_{\vb{p}}E_{\vb{p}}
 \hat a_{\vb{p}}^\dagger\hat a_{\vb{p}}.}
 \label{eq:scalar-H-excitation-dp76}
\end{equation}
在固定几何下，正规序把共同的真空常数与粒子激发能分开：将湮灭算符移到产生算符右侧后，自由标量哈密顿量的激发部分正是\eqnrefs{eq:scalar-H-excitation-dp76}。改变边界、几何或引力耦合时，真空能差需要重新作为物理量处理。

空间平移的 Noether 荷量子化后，动量算符为
\begin{equation}
 \hat{\vb{P}}
 =\int\frac{\dd^3 p}{(2\pi\hbar)^3}
 \vb{p}\,\hat a^\dagger(\vb{p})\hat a(\vb{p}).
 \label{eq:scalar-P-r3}
\end{equation}
于是
\begin{equation*}
 |\vb{p}\rangle\equiv\hat a^\dagger(\vb{p})|0\rangle
\end{equation*}
同时是能量与动量的单量子本征态，并满足
\begin{equation*}
 \hat H_0^{\rm exc}|\vb{p}\rangle=E_{\vb{p}}|\vb{p}\rangle,
 \qquad
 \hat{\vb{P}}|\vb{p}\rangle=\vb{p}|\vb{p}\rangle.
\end{equation*}
因此“粒子”不是量子化以前额外加入的对象，而是自由场时空平移生成元的共同本征激发。对标量场，静止单量子态的内禀角动量为零，所以这些激发具有自旋 $0$。

\subsection{Fock 空间与态归一化}

场量子化把每个正常模变成谐振子，但散射理论真正使用的是由这些模式构成的多粒子 Hilbert 空间（希尔伯特空间）。定义真空 $a(\vb{p})|0\rangle=0$，一粒子态由 $a^\dagger(\vb{p})|0\rangle$ 产生。实标量场的产生算符彼此对易，因此 $n$ 粒子扇区是单粒子 Hilbert 空间的对称张量幂：
\begin{equation*}
\mathscr F_B
=\mathbb C|0\rangle\oplus\mathcal H_1
\oplus\operatorname{Sym}\mathcal H_1^{\otimes2}\oplus\cdots .
\end{equation*}
自由理论中
\begin{equation*}
N=\int\frac{\dd^3 p}{(2\pi\hbar)^3}a^\dagger(\vb{p})a(\vb{p})
\end{equation*}
与 $H_0$ 对易；相互作用则一般混合不同粒子数扇区，因此 Fock 空间而不是固定 $N$ 的 Hilbert 空间才是相对论场论的自然状态空间。

自由哈密顿量含有每个模式的零点能。场论计算中经常需要一种统一的算符排列规则：把所有产生算符移动到所有湮灭算符左侧；若交换的是费米算符，每交换一次还要带上反对易产生的负号。这种排列称为\emph{正规序}，并用一对冒号 $:\!\cdots\!:$ 表示。对当前自由标量哈密顿量，正规序恰好把由 $aa^\dagger=a^\dagger a+1$ 产生的真空 $c$ 数常数与激发部分分开，因此它与正文已经定义的激发哈密顿量完全相同：
\begin{equation}
{
:\!H_0\!:
=\int\frac{\dd^3 p}{(2\pi\hbar)^3}
E_{\vb{p}}\,a^\dagger(\vb{p})a(\vb{p}),
\qquad \langle0|:\!H_0\!:|0\rangle=0.}
\label{eq:normal-H-r2}
\end{equation}
正规序选定了算符排列与参考能量零点；它并不消除真空能在所有物理问题中的作用。改变边界条件会改变模谱，两个真空能之间的有限差仍可产生 Casimir 力；进入有引力或重整化问题时，真空能本身也需要作为理论参数处理。

连续动量态是广义本征态，而不是单位范数 Hilbert 空间向量。Lorentz 对称性要求散射完备关系使用正能质量壳上的不变测度
\begin{equation}
{
\frac{\dd^3 p}{(2\pi\hbar)^3}\frac{c}{2E_{\vb{p}}}.}
\label{eq:lorentz-measure-r2}
\end{equation}
相应一粒子归一化与完备关系为
\begin{equation*}
\langle p|q\rangle
=(2\pi\hbar)^3\frac{2E_{\vb{p}}}{c}
\delta^{(3)}(\vb{p}-\vb{q}),
\end{equation*}
以及


\begin{proposition}[连续一粒子态的完备测度]
设连续动量态按正文的 Lorentz 不变内积归一化；完备关系在一粒子 Hilbert 空间及其稠密检验域上理解。则有
\begin{equation}
{
\id(1)
=\int\frac{\dd^3 p}{(2\pi\hbar)^3}\frac{c}{2E_{\vb{p}}}
|p\rangle\langle p|.}
\label{eq:one-particle-completeness-r2}
\end{equation}
\end{proposition}
\begin{proof}
正文\eqnrefs{eq:lorentz-measure-r2} 可以从四维 Lorentz 不变量
\[
\dd^4 p\,\delta(p^2-m^2c^2)\Theta(p^0)
\]
直接读出。固定 $\vb{p}$ 后，质量壳方程在 $p^0$ 上有两个根
\[
p^0=\pm\frac{E_{\vb{p}}}{c},
\qquad
E_{\vb{p}}=\sqrt{m^2c^4+c^2\vb{p}^2}.
\]
利用
\[
\delta[f(x)]=\sum_i\frac{\delta(x-x_i)}{|f'(x_i)|}
\]
以及 $f(p^0)=(p^0)^2-E_{\vb{p}}^2/c^2$，有
\begin{align*}
\delta\!\left((p^0)^2-\frac{E_{\vb{p}}^2}{c^2}\right)
&=\frac{c}{2E_{\vb{p}}}
\left[
\delta\!\left(p^0-\frac{E_{\vb{p}}}c\right)
+\delta\!\left(p^0+\frac{E_{\vb{p}}}c\right)
\right].
\end{align*}
$\Theta(p^0)$ 只保留正能根，因此 $p^0$ 积分给出正文测度中的 $c/(2E_{\vb{p}})$。

现在考虑正文已经给出的协变归一化
\[
\langle p|q\rangle
=(2\pi\hbar)^3\frac{2E_{\vb{p}}}{c}
\delta^{(3)}(\vb{p}-\vb{q}).
\]
令
\[
\mathcal I_1
\equiv
\int\frac{\dd^3 p}{(2\pi\hbar)^3}
\frac{c}{2E_{\vb{p}}}|p\rangle\langle p|.
\]
让它作用在任意基矢 $|q\rangle$ 上：
\begin{align*}
\mathcal I_1|q\rangle
&=\int\frac{\dd^3 p}{(2\pi\hbar)^3}
\frac{c}{2E_{\vb{p}}}|p\rangle
(2\pi\hbar)^3\frac{2E_{\vb{p}}}c
\delta^{(3)}(\vb{p}-\vb{q})\\
&=\int\dd^3 p\,|p\rangle
\delta^{(3)}(\vb{p}-\vb{q})\\
&=|q\rangle.
\end{align*}
由于这些正能质量壳基矢张成一粒子扇区，$\mathcal I_1$ 就是一粒子恒等算符，从而得到正文\eqnrefs{eq:one-particle-completeness-r2}。
\end{proof}

把场的相关函数转换成散射振幅时，外部单粒子从渐近区传播到相互作用区的自由极点因子必须被约去。这个操作称为外腿约化，并且必须与 Lorentz 不变相空间和截面使用同一套一粒子态归一化。因此态归一化不是记号细节，而是所有 $2E$ 与 $(2\pi\hbar)^3$ 因子的共同来源。

最后，局域场 $\varphi(x)$ 严格说是算符值分布。物理探测器具有有限时空分辨率，因此真正良定义的对象是经测试函数涂抹的场
\begin{equation*}
\varphi[f]=\int \dd^4 x\,f(x)\varphi(x),
\end{equation*}
这也解释了为什么固定点的 $\langle\varphi(x)^2\rangle$ 发散并不直接等同于可测量发散。

\begin{exbox}[一维 Casimir 模型]
零点能为何不能简单“全部删除”可由一维 Dirichlet 标量模型看出。间距为 $d$ 的两边界允许 $k_n=n\pi/d$，形式零点能为
\begin{equation*}
E_0(d)=\frac{\pi\hbar c}{2d}\sum_{n=1}^\infty n.
\end{equation*}
求两边界之间的真空力，并区分调节器相关的真空能与有限的边界力。
\textbf{解。} 引入长度调节器 $a>0$，
\begin{equation*}
E_0(d;a)=\frac{\pi\hbar c}{2d}
\sum_{n=1}^\infty n \ee^{-a\pi n/d}
=\frac{\pi\hbar c}{2d}
\frac{\ee^{-a\pi/d}}{(1-\ee^{-a\pi/d})^2}.
\end{equation*}
小 $a$ 展开给
\begin{equation*}
E_0(d;a)=\frac{\hbar c\,d}{2\pi a^2}-\frac{\pi\hbar c}{24d}+O(a^2/d^3).
\end{equation*}
把系统嵌入总长 $L$ 的大盒后，$d$ 与 $L-d$ 两区的 $1/a^2$ 项只合成与 $d$ 无关的常数；对 $d$ 求导并取 $L\to\infty,a\to0$，得到
\begin{equation*}
F(d)=-\frac{\partial E}{\partial d}=-\frac{\pi\hbar c}{24d^2}.
\end{equation*}
因此，依赖调节器的绝对真空能量与不依赖调节器的边界力可以被明确区分。
\end{exbox}

实标量 Fock 空间还留下一个自然问题：怎样在同一套量子化语言里表示带守恒电荷的粒子与反粒子？答案不是另造一套状态空间，而是把两个实自由度组合成一个复场。

实标量场只有一个实自由度，不能在保持实条件的同时承受非平凡的连续相位旋转；因此它没有像电荷那样的 $U(1)$ Noether 荷。若希望最简单的自旋零场也能携带守恒内部荷，就必须把两个实场组合成一个复场。这样正、负频系数不再互为同一个算符的厄米共轭，而会形成两套独立的粒子与反粒子算符。

实标量场满足 $\varphi^\dagger=\varphi$，正负频系数因此互为厄米共轭。复标量场 $\varphi_{\rm ch}$ 不受这个条件约束。把它的静质量记为 $m_{\rm ch}$，其中 $[m_{\rm ch}]=\mathrm{kg}$；这个记号随后也用于复标量玩具模型。取
\begin{equation*}
 \Lag_{\rm ch}
 =\partial_\mu\varphi_{\rm ch}^*\partial^\mu\varphi_{\rm ch}
 -\frac{m_{\rm ch}^2c^2}{\hbar^2}\varphi_{\rm ch}^*\varphi_{\rm ch},
\end{equation*}
其量子场展开需要两套独立算符：
\begin{equation}
 \hat\varphi_{\rm ch}(x)=
 \int\frac{\dd^3 p}{(2\pi\hbar)^3}
 \frac{\hbar c}{\sqrt{2E_{\vb{p}}}}
 \left[
 \hat b(\vb{p})\ee^{-\ii p\cdot x/\hbar}
 +\hat d^\dagger(\vb{p})\ee^{+\ii p\cdot x/\hbar}
 \right].
 \label{eq:complex-scalar-exp-r2}
\end{equation}
$\hat b^\dagger$ 与 $\hat d^\dagger$ 分别产生粒子和反粒子。全局 $U(1)$ 变换 $\varphi_{\rm ch}\to \ee^{\ii\alpha}\varphi_{\rm ch}$ 的 Noether 电荷量子化后具有形式
\begin{equation}
 \hat Q=q\int\frac{\dd^3 p}{(2\pi\hbar)^3}
 \left[
 \hat b^\dagger\hat b-\hat d^\dagger\hat d
 \right].
 \label{eq:complex-scalar-charge-r2}
\end{equation}
因此粒子数 $N_b$ 和反粒子数 $N_d$ 在有相互作用时可以分别改变，而净 $U(1)$ 电荷仍可守恒。反粒子并不是 Dirac 海或 Fermi 统计独有的现象；它已经出现在最简单的相对论复标量场中。

\begin{exbox}[粒子与反粒子的 $U(1)$ 荷]
单看\eqnrefs{eq:complex-scalar-charge-r2} 中的减号还不够；真正可测的是荷算符作用在激发态上的本征值。定义
\begin{equation*}
 |\vb{p};b\rangle
 =\hat b^\dagger(\vb{p})|0\rangle,
 \qquad
 |\vb{p};d\rangle
 =\hat d^\dagger(\vb{p})|0\rangle,
 \qquad \hat Q|0\rangle=0.
\end{equation*}
求两种单粒子激发的荷本征值，说明反粒子的电荷符号与能量符号怎样区别。
\textbf{解。} 由玻色对易关系，数算符满足
\begin{align*}
 [\hat b^\dagger\hat b,\hat b^\dagger]&=+\hat b^\dagger,\\
 [\hat d^\dagger\hat d,\hat d^\dagger]&=+\hat d^\dagger.
\end{align*}
因此荷算符与两种产生算符的对易子分别为
\begin{equation*}
 [\hat Q,\hat b^\dagger(\vb{p})]
 =+q\,\hat b^\dagger(\vb{p}),
 \qquad
 [\hat Q,\hat d^\dagger(\vb{p})]
 =-q\,\hat d^\dagger(\vb{p}).
\end{equation*}
把 $\hat Q$ 移过产生算符作用到真空上，便得到
\begin{equation*}
 \hat Q|\vb{p};b\rangle=+q|\vb{p};b\rangle,
 \qquad
 \hat Q|\vb{p};d\rangle=-q|\vb{p};d\rangle.
\end{equation*}
这个结果把负频率分支与反粒子携带相反内部荷区分开：反粒子的能量仍由 $E_{\vb{p}}>0$ 给出，改变符号的是守恒荷本征值。Dirac 场中的正电子遵循同一粒子--反粒子逻辑，只是算符代数改为费米反对易关系。
\end{exbox}

平面波量子化已经固定了等时对易关系与能量归一化。截面改变时，这些条件仍须保持；改变的是空间算子的谱及定义域。以下采用与原标量场相容的规范化 $\varphi=\phi_{\rm phys}/\sqrt{\hbar c}$，先在有限模态中建立一般截面的振子，再按谱测度处理连续谱。

\subsection{由正自伴空间算子量子化一般截面}
\label{subsec:ghs-mode-quantization}


先取一个静态的 $(n+1)$ 维洛伦兹时空 $\RR\times\Sigma$，其中时间坐标是具有长度量纲的 $x^0=ct$，$\Sigma$ 是一个紧致、无边界的 $n$ 维黎曼空间，度规为 $h_{ab}$。这里 $n$ 是空间维数，后文讨论欧氏共形场论时则以 $d$ 表示整个欧氏空间维数；在两个问题之间转换时，不应提前把 $n$ 和 $d$ 混用。为使谱从正数开始，先考虑严格正的空间算子
\begin{equation}\label{eq:ghs-bridge-spatial-A}
 A=-\Delta_h+\mu^2,\qquad \mu>0,
\end{equation}
其中 $\mu$ 具有逆长度量纲；若场的物理质量为 $m$，则 $\mu=mc/\hbar$。$\Delta_h=\operatorname{div}_h\nabla_h$ 的本征值取非正的号。对应的无量纲作用量（在合适的场归一化后）为
\begin{equation}\label{eq:ghs-bridge-lorentz-action}
 \mathscr S[\varphi]=\frac12\int\dd x^0\,\dd\mu_h
 \bigl[(\partial_0\varphi)^2-h^{ab}(\partial_a\varphi)(\partial_b\varphi)-\mu^2\varphi^2\bigr].
\end{equation}
物理作用量为 $S=\hbar\mathscr S$；规范化场 $[\varphi]=L^{(1-n)/2}$，因而 $[q_j]=L^{1/2}$、$[p_j]=L^{-1/2}$，$[q_j,p_k]=\ii\delta_{jk}$ 的两边均无量纲。这里 $\dd\mu_h=\sqrt{\det h}\,\dd^ny$。本节选择规范化振子变量；恢复物理量时每个振子的能量是 $\hbar c\omega$，其中 $\omega$ 是逆长度本征频率。以 $x^0$ 作为演化变量使 $c$ 与 $\hbar$ 不会在谱的定义中被暗中设为一。

在无边界紧空间上，$-\Delta_h$ 的非负本征值是离散的。记 $u_j\in L^2(\Sigma,\dd\mu_h)$ 为一组实正交归一的本征函数，重数计入指标 $j$，满足
\begin{equation}\label{eq:ghs-bridge-eigenbasis}
 Au_j=\omega_j^2u_j,\quad \omega_j>0,\qquad
 \int_\Sigma u_j u_k\,\dd\mu_h=\delta_{jk},\qquad
 \sum_j u_j(y)u_j(y')=\delta_h(y,y').
\end{equation}
最后一条是分布意义的完备性关系：$\int\delta_h(y,y')f(y')\dd\mu_h(y')=f(y)$。$u_j$ 是空间函数，$\omega_j$ 是频率而非能量；二者已经足以代替平直空间的 Fourier 波矢。

将场展开为 $\varphi(x^0,y)=\sum_jq_j(x^0)u_j(y)$。同一展开给出
\begin{align*}
 \int_\Sigma(\partial_0\varphi)^2\dd\mu_h
 &=\sum_{j,k}\dot q_j\dot q_k\int_\Sigma u_ju_k\dd\mu_h
 =\sum_j\dot q_j^{\,2},\\
 \int_\Sigma\bigl(|\nabla\varphi|_h^2+\mu^2\varphi^2\bigr)\dd\mu_h
 &=\int_\Sigma\varphi A\varphi\dd\mu_h\\
 &=\sum_{j,k}q_jq_k\omega_k^2\int_\Sigma u_ju_k\dd\mu_h
 =\sum_j\omega_j^2q_j^2.
\end{align*}
第二行的积分分部在无边界紧空间上没有边界项；当 $\Sigma$ 有边界时，必须先选择使 $A$ 自伴的条件，不能直接使用这一等式。于是场作用量完全化成
\begin{equation}\label{eq:ghs-bridge-oscillator-action}
 \mathscr S=\frac12\int\dd x^0\sum_j(\dot q_j^2-\omega_j^2q_j^2).
\end{equation}
这就是“自由场是无限多个振子”的一般数学含义：真正重要的条件是空间算子存在合适的谱分解，而非平面波是否存在。



为避免把无限维问题直接塞进未经定义的乘积，先只保留前 $N$ 个模。共轭动量为 $p_j=\dot q_j$，以逆长度计的哈密顿量为 $H_N=\frac12\sum_{j=1}^{N}(p_j^2+\omega_j^2q_j^2)$。施加 $[q_j,p_k]=\ii\delta_{jk}$，定义
\begin{align*}
 a_j&=\sqrt{\frac{\omega_j}{2}}q_j+\frac{\ii}{\sqrt{2\omega_j}}p_j,\\
 a_j^\dagger&=\sqrt{\frac{\omega_j}{2}}q_j-\frac{\ii}{\sqrt{2\omega_j}}p_j.
\end{align*}
逐项计算交叉对易子得
\begin{align*}
 [a_j,a_k^\dagger]
 &=-\frac{\ii}{2}\sqrt{\frac{\omega_j}{\omega_k}}[q_j,p_k]
 +\frac{\ii}{2}\sqrt{\frac{\omega_k}{\omega_j}}[p_j,q_k]
 =\delta_{jk}.
\end{align*}
反解为 $q_j=(a_j+a_j^\dagger)/\sqrt{2\omega_j}$ 和
$p_j=-\ii\sqrt{\omega_j/2}(a_j-a_j^\dagger)$，代回 $H_N$ 可得
\[
 H_N=\sum_{j=1}^{N}\omega_j\left(a_j^\dagger a_j+\frac12\right).
\]
物理哈密顿量为 $\hbar cH_N$。零点项 $\frac12\sum_j\omega_j$ 在 $N\to\infty$ 时通常发散；后文讨论的正规序哈密顿量只计算相对真空的激发能量，不等于已经解决了一般背景上的真空能量重整化。



有限模真空 $|0_N\rangle$ 由 $a_j|0_N\rangle=0$ 定义。由此制备的任意有限粒子态可以通过产生算符作用于真空得到。无限模极限不再只是形式上的 $N\to\infty$：要得到真正的 Fock 表示，还必须讨论哪些序列有有限范数、不同正频分解之间能否由酉算子联系。这正是后文复结构与 Bogoliubov 理论要回答的问题。


要理解欧氏径向量子化，先看一个振子的真空在坐标表象中是什么。令 $p_j=-\ii\partial_{q_j}$。条件 $a_j\psi_0(q)=0$ 就是
\[
 \left(\sqrt{\frac{\omega_j}{2}}q_j+\frac1{\sqrt{2\omega_j}}\partial_{q_j}\right)\psi_0(q_j)=0,
 \qquad \partial_{q_j}\log\psi_0=-\omega_jq_j.
\]
积分后，利用 $\int_\RR\ee^{-\omega q^2}\dd q=\sqrt{\pi/\omega}$ 归一化，得到
\begin{equation}\label{eq:ghs-bridge-vacuum-one-mode}
 \psi_{0,j}(q_j)=\left(\frac{\omega_j}{\pi}\right)^{1/4}\ee^{-\omega_jq_j^2/2}.
\end{equation}
$N$ 个独立模的真空波函数是这些高斯函数的乘积，其指数为 $-\frac12\sum_{j=1}^{N}\omega_jq_j^2$。将 $q_j=\langle u_j,f\rangle_{L^2}$ 视为边界场 $f(y)$ 的谱系数，就得到形式表达
\begin{equation}\label{eq:ghs-bridge-vacuum-functional}
 \Psi_0[f]\ \propto\
 \exp\left[-\frac12\langle f,A^{1/2}f\rangle_{L^2(\Sigma)}\right].
\end{equation}
这里的 $A^{1/2}$ 定义为 $A^{1/2}u_j=\omega_ju_j$。它一般不是局域微分算子：$A$ 是二阶算子，而其正平方根通常在一点上的取值依赖整个截面上的 $f$。因此真空波泛函的高斯核已经含有整个空间的几何信息。



式 \cref{eq:ghs-bridge-vacuum-functional} 首先是有限模正则化下的高斯泛函。无限维空间没有平移不变的 Lebesgue 测度 $\mathcal Df$；严谨构造应通过圆柱高斯测度或相应的算子代数完成。若 $A^{-1/2}$ 不是迹类算子，典型的真空场不能视为普通 $L^2$ 函数，而是分布。后文的谱核在分离点可按分布方式解释。



波泛函的指数为什么恰好出现平方根，而不是 $A$ 本身？因为 $A$ 给出振子的\emph{频率平方}，但高斯真空由一阶方程 $a_j\psi_0=0$ 决定。取平方根正是把二阶场方程分成正频与负频两部分的一次运算。稍后在一般曲时空上，若不存在一个全局不变的正频划分，真空也就不再有这个唯一的形式。


在初等场论中，我们常在动量空间写传播子 $(p^2+m^2)^{-1}$。如果空间没有平移对称性，动量就不是好量子数；然而“传播子是微分算子的逆”仍然成立。先作 Wick 转动，将 $x^0$ 换为欧氏时间 $\tau$，二次作用量是
\begin{equation}\label{eq:ghs-bridge-Euclidean-action}
 \mathscr I_E=\frac12\int_{\RR}\dd\tau\int_\Sigma\dd\mu_h
 \left[(\partial_\tau\varphi)^2+\varphi A\varphi\right].
\end{equation}
对应欧氏算子为 $L=-\partial_\tau^2+A$。我们要找的两点函数 $G$ 应满足
\begin{equation}\label{eq:ghs-bridge-green-equation}
 L_{(\tau,y)}G(\tau,y;\tau',y')
 =\delta(\tau-\tau')\,\delta_h(y,y').
\end{equation}
这条定义明确了逆算子，也说明传播子依赖边界条件：在整条欧氏时间轴上取两端衰减的解，与有限柱体上固定边界场的解并不相同。

将 $G$ 沿空间本征基展开为 $G=\sum_jg_j(\tau,\tau')u_j(y)u_j(y')$，代入 \cref{eq:ghs-bridge-green-equation} 并使用正交性，可逐模得到
\[
 (-\partial_\tau^2+\omega_j^2)g_j(\tau,\tau')=\delta(\tau-\tau').
\]
当 $\tau\ne\tau'$ 时，齐次方程只允许 $\ee^{\pm\omega_j\tau}$。要求 $|\tau|\to\infty$ 衰减，则解应与 $\ee^{-\omega_j|\tau-\tau'|}$ 成正比。记系数为 $c_j$，将方程跨越 $\tau'$ 的一个小区间积分，得到导数跳跃条件
\[
 -\left[\partial_\tau g_j\right]_{\tau'^-}^{\tau'^+}=1,
 \qquad 2\omega_jc_j=1.
\]
因此任意严格正自伴空间算子的欧氏圆柱传播核为
\begin{equation}\label{eq:ghs-bridge-spectral-G}
 G(\tau,y;\tau',y')
 =\sum_j\frac{\ee^{-\omega_j|\tau-\tau'|}}{2\omega_j}
 u_j(y)u_j(y').
\end{equation}
这就是熟悉的动量积分在一般紧致空间上的替代物。若空间非紧，谱可能包含连续部分，此时必须将求和替换为投影谱测度的积分，不能人为把连续谱离散化。


这一离散本征函数展开还可改写为不依赖球面坐标的谱积分。
对正自伴算子 $A$，谱定理给出投影值测度 $E_A(\lambda)$，并定义 $F(A)=\int_{(0,\infty)}F(\lambda)\,\dd E_A(\lambda)$。取 $F(\lambda)=\ee^{-|t|\sqrt\lambda}/(2\sqrt\lambda)$，便得到
\[
 G(t)=\frac12A^{-1/2}\ee^{-|t|A^{1/2}}
 =\int_{(0,\infty)}\frac{\ee^{-|t|\sqrt\lambda}}{2\sqrt\lambda}\,\dd E_A(\lambda).
\]
当 $A\ge\mu^2I>0$ 时，$A^{-1/2}$ 有界。对 $t\ne0$，函数演算允许逐谱求导：
\[
 (\partial_t^2-A)G(t)=0.
\]
而一侧导数为 $G'(0^+)=-\frac12I$、$G'(0^-)={+}\frac12I$，故分布二阶导数含 $-\delta(t)I$，从而 $(-\partial_t^2+A)G=\delta(t)I$。


这一计算展示了谱理论在场论中的真正用途：当不能求出所有本征函数时，依然可以用函数演算定义传播子；当可以显式求出谱时，又能够立即得到可计算的级数或积分。谱定理所涉及的定义域、自伴性和连续谱可参照《数学技巧》中自伴算子与谱测度的完整论证。




电子的非相对论动力学已经由 Dirac 哈密顿量的 Foldy--Wouthuysen 分块对角化得到，最低阶即为正能 Pauli--Schr\"odinger 哈密顿量。此处要连接的是态空间：普通的一粒子波函数如何作为 Fock 空间固定粒子数扇区的坐标表示出现。守粒子数的非相对论场给出这个对应。

% -----------------------------------------------------------------------------
% Beginner bridge: how ordinary one-particle quantum mechanics emerges
% -----------------------------------------------------------------------------

量子场算符可以改变粒子数，而普通量子力学的波函数只是复数值函数；两种语言之间的联系最好先在一个守粒子数的非相对论 Schr\"odinger 场中看清。这个低能模型把 Fock 空间按粒子数分解，并使普通单粒子量子力学直接作为其中的固定一粒子扇区出现。

等时正则关系为
\begin{equation}
 [\hat\Psi(\vb{r}),\hat\Psi^\dagger(\vb{r}')]
 =\delta^{(3)}(\vb{r}-\vb{r}'),
 \qquad
 [\hat\Psi,\hat\Psi]=0.
 \label{eq:sch-field-comm-r6}
\end{equation}
若真空满足 $\hat\Psi(\vb{r})|0\rangle=0$，则 $\hat\Psi^\dagger(\vb{r})|0\rangle$ 是在位置标签 $\vb{r}$ 上产生一个粒子的广义态。任意一粒子态因此可展开为
\begin{equation}
 {
 |\Psi_1(t)\rangle
 =\int\dd^3 r\,
 \psi_{\rm 1p}(\vb{r},t)
 \hat\Psi^\dagger(\vb{r})|0\rangle .}
 \label{eq:fock-oneparticle-expansion-dp68}
\end{equation}
这里的 $\psi_{\rm 1p}$ 才是普通第一量子化中的复数值波函数；$\hat\Psi$ 本身仍是作用在 Fock 空间上的算符值分布。

对自由哈密顿量
\begin{equation*}
 \hat H
 =\int\dd^3 r\,
 \hat\Psi^\dagger(\vb{r})
 \left(-\frac{\hbar^2\nabla^2}{2m}\right)
 \hat\Psi(\vb{r}),
\end{equation*}
Schr\"odinger 绘景的态方程在一粒子扇区中约化为


\begin{proposition}[二次量子化在一粒子空间上的限制]
设一粒子算子为自由非相对论哈密顿量 $h=-\hbar^2\nabla^2/(2m)$，取本节指定边界条件下的自伴定义域，并令 $\hat H=\dd\Gamma(h)$。一粒子态及其作用均属于相应定义域，下式成立；一般自伴 $h$ 的对应一粒子方程为 $\ii\hbar\partial_t\psi=h\psi$：
\begin{equation}
 {
 \ii\hbar\partial_t\psi_{\rm 1p}(\vb{r},t)
 =-\frac{\hbar^2}{2m}\nabla^2\psi_{\rm 1p}(\vb{r},t).}
\label{eq:fock-oneparticle-schrodinger-dp68}
\end{equation}
\end{proposition}
\begin{proof}
一粒子波函数不是额外附加到场论上的对象，而是 Fock 态在局域一粒子基底上的坐标。先从正文的态展开
\begin{equation*}
|\Psi_1(t)\rangle
=\int\dd^3r\,\psi_{\rm 1p}(\vb{r},t)
\hat\Psi^\dagger(\vb{r})|0\rangle
\end{equation*}
出发。自由 Schr\"odinger 场哈密顿量为
\begin{equation*}
\hat H
=\int\dd^3x\,
\hat\Psi^\dagger(\vb{x})
\left(-\frac{\hbar^2\nabla_{\vb{x}}^2}{2m}\right)
\hat\Psi(\vb{x}).
\end{equation*}
利用等时对易关系，先求产生一个局域粒子时能量算符怎样作用：
\begin{align*}
[\hat H,\hat\Psi^\dagger(\vb{r})]
&=\int\dd^3x\,
\hat\Psi^\dagger(\vb{x})
\left(-\frac{\hbar^2\nabla_{\vb{x}}^2}{2m}\right)
[\hat\Psi(\vb{x}),\hat\Psi^\dagger(\vb{r})]\\
&=\int\dd^3x\,
\hat\Psi^\dagger(\vb{x})
\left(-\frac{\hbar^2\nabla_{\vb{x}}^2}{2m}\right)
\delta^{(3)}(\vb{x}-\vb{r})\\
&=-\frac{\hbar^2}{2m}\nabla_{\vb{r}}^2
\hat\Psi^\dagger(\vb{r}).
\end{align*}
因为真空满足 \(\hat H|0\rangle=0\)，于是
\begin{align*}
\hat H|\Psi_1\rangle
&=\int\dd^3r\,\psi_{\rm 1p}(\vb{r})
[\hat H,\hat\Psi^\dagger(\vb{r})]|0\rangle\\
&=\int\dd^3r\,
\left[-\frac{\hbar^2}{2m}\nabla^2\psi_{\rm 1p}(\vb{r})\right]
\hat\Psi^\dagger(\vb{r})|0\rangle,
\end{align*}
其中第二行对空间变量分部积分两次，并假设波函数在边界处足够快衰减。

连续基底中的系数比较由场算符投影严格实现。对任意 \(\vb{x}\)，左乘 \(\langle0|\hat\Psi(\vb{x})\)，由
\begin{equation*}
\langle0|\hat\Psi(\vb{x})
\hat\Psi^\dagger(\vb{r})|0\rangle
=\delta^{(3)}(\vb{x}-\vb{r})
\end{equation*}
可得
\begin{align*}
\langle0|\hat\Psi(\vb{x})
\ii\hbar\partial_t|\Psi_1\rangle
&=\ii\hbar\partial_t\psi_{\rm 1p}(\vb{x},t),\\
\langle0|\hat\Psi(\vb{x})
\hat H|\Psi_1\rangle
&=-\frac{\hbar^2}{2m}\nabla_{\vb{x}}^2
\psi_{\rm 1p}(\vb{x},t).
\end{align*}
因此 Fock 空间 Schr\"odinger 方程 \(\ii\hbar\partial_t|\Psi_1\rangle=\hat H|\Psi_1\rangle\) 在每个 \(\vb{x}\) 上严格投影成
\begin{equation*}
\ii\hbar\partial_t\psi_{\rm 1p}(\vb{x},t)
=-\frac{\hbar^2}{2m}\nabla_{\vb{x}}^2\psi_{\rm 1p}(\vb{x},t),
\end{equation*}
即正文\eqnrefs{eq:fock-oneparticle-schrodinger-dp68}。这里真正建立的是“固定粒子数 Fock 扇区”与普通单粒子波函数之间的坐标同构，而不是重新假设一条量子力学方程。
\end{proof}

\eqnrefs{eq:fock-oneparticle-schrodinger-dp68} 不是另加的一条动力学假设，而是把场论哈密顿量限制到\eqnrefs{eq:fock-oneparticle-expansion-dp68} 所定义的一粒子扇区后的结果。

$\hat\Psi^\dagger$ 用来构造粒子态，$\psi_{\rm 1p}$ 用来描述固定一粒子态在位置基中的展开系数。进入两粒子及更高扇区后，态由多个产生算符构造，玻色对易关系或费米反对易关系自动决定交换对称性，因此不需要再把场算符误解成一条“多粒子波函数”。


固定粒子数扇区解释了普通波函数从哪里来。下一步回到相对论场论：此时粒子数不再必须守恒，真正需要保留的是局域场及其二点函数；类空分离处的场对易关系将给出微观因果性。


Fourier 变换、分布、$\ii0^+$ 与围道规定由数学母定义固定；自由场传播子由这些数学边界条件作用于相应逆算符而得到。

传播子属于时空二点函数，所以先转到 Heisenberg 绘景。Schr\"odinger 绘景的场算符在固定时刻定义，而 Heisenberg 场定义为
\begin{equation*}
 \hat\varphi_H(\vb{x},t)
 =\ee^{+\ii\hat H_0t/\hbar}\hat\varphi_S(\vb{x})
 \ee^{-\ii\hat H_0t/\hbar}.
\end{equation*}
由 $[\hat H_0,\hat a(\vb{p})]=-E_{\vb{p}}\hat a(\vb{p})$ 得
\begin{equation*}
 \ee^{+\ii\hat H_0t/\hbar}\hat a(\vb{p})\ee^{-\ii\hat H_0t/\hbar}
 =\ee^{-\ii E_{\vb{p}}t/\hbar}\hat a(\vb{p}),
\end{equation*}
于是恢复了 \eqnrefs{eq:scalar-mode-r2} 的完整时空平面波。进一步计算任意两点的场对易子
\begin{align}
 [\hat\varphi(x),\hat\varphi(y)]
 &=\hbar^2c^2\int\frac{\dd^3 p}{(2\pi\hbar)^3}
 \frac{1}{2E_{\vb{p}}}
 \left[\ee^{-\ii p\cdot(x-y)/\hbar}-\ee^{+\ii p\cdot(x-y)/\hbar}\right]\\
 &\equiv \ii\hbar c\,\Delta(x-y).
 \label{eq:pauli-jordan-r3}
\end{align}
积分测度 $\dd^3 p/(2E_{\vb{p}})$ 是 Lorentz 不变，而在等时、空间分离情况下，把第二项作 $\vb{p}\to-\vb{p}$ 变换便与第一项抵消，所以
\begin{equation}
 [\hat\varphi(x),\hat\varphi(y)]=0,
 \qquad (x-y)^2<0.
 \label{eq:microcausality-scalar-r3}
\end{equation}
任何类空分离都能通过 Lorentz 变换化到等时分离，因此上式对全部类空点成立。这是局域量子场论的微观因果性条件：类空分离处的局域玻色可观测量可以同时测量而不彼此影响。

现在定义不带时间排序的真空二点函数
\begin{equation}
 D^+(x-y)\equiv
 \langle0|\hat\varphi(x)\hat\varphi(y)|0\rangle
 =\hbar^2c^2\int\frac{\dd^3 p}{(2\pi\hbar)^3}
 \frac{\ee^{-\ii p\cdot(x-y)/\hbar}}{2E_{\vb{p}}}.
 \label{eq:wightman-plus-r3}
\end{equation}
它常称为正频 Wightman 函数。它可以在类空分离时非零，因此不能被解释成“粒子以超光速从 $y$ 跑到 $x$ 的概率幅”。真正控制因果性的对象是两种顺序之差
\begin{equation*}
 D^+(x-y)-D^+(y-x)
 =\langle0|[\hat\varphi(x),\hat\varphi(y)]|0\rangle,
\end{equation*}
其类空值由 \eqnrefs{eq:microcausality-scalar-r3} 为零。Feynman 传播子把两个 Wightman 函数按时间先后拼接，从而同时编码传播边界条件和时序乘积。

% Detailed derivation gathered at the end of the owning structure.


